\section{从东京出发}

\begin{event}{\eventref{NS-0960-CHENQIAO}{陈桥建宋}（九六〇年正月；后周禁军与赵匡胤集团；陈桥、开封；建国可核，军中密议与动机未能复原）}
《宋史·太祖本纪》《续资治通鉴长编》建隆元年正月条及《宋会要辑稿·礼》可以交叉追踪即位、年序与礼仪语言。它们成书层次不同；“黄袍加身”的仪式细节和全体参与者的意图，不能据后出的铺陈写成现场。
\end{event}

公元960年正月，赵匡胤在陈桥兵变后即位，国号宋。\hyperlink{evt:NS-0960-CHENQIAO}{陈桥建宋（960）}新朝首先占有的，并不只是一顶皇帝的冠冕，而是开封这座已经运转的都城：宫城、官署、仓储、汴河和驻守其间的禁军，把政令、粮食与兵员系在同一处。东京城遗址所见的城垣、街道和水系，可以帮助理解这种承继；它并不容许我们把北宋都城的每一层空间都写得仿佛亲见。宋廷承用开封为东京，正因为这里既是后周的中枢，也是维持百官与军队供给的节点。

\phantomsection\label{fig:vol04:northern-song:founding-chenqiao-kaifeng}
\begin{center}
\begin{tikzpicture}[x=.88cm,y=.88cm,font=\small]
  \fill[paper] (0,0) rectangle (11.6,6.25);

  \fill[source!16] (.55,3.85) -- (3.65,4.15) -- (4.1,5.75) -- (.8,5.75) -- cycle;
  \node[align=center] at (2.25,5.0) {北方行军地带\\后周禁军北伐动员};
  \fill[timeline!23] (.75,.7) -- (4.35,.85) -- (4.1,3.15) -- (.55,3.0) -- cycle;
  \node[align=center] at (2.35,1.9) {陈桥\\军队驻泊与政权转换};
  \fill[dynasty!24] (7.1,1.0) -- (10.95,1.2) -- (11.05,5.35) -- (7.45,5.2) -- cycle;
  \node[align=center] at (9.1,3.65) {开封（东京）\\宫城、官署、仓储\\与禁军供给节点};

  \draw[very thick,color=source!75] (4.2,.6) .. controls (5.7,1.8) and (6.2,4.45) .. (7.1,5.85);
  \node[text=source,rotate=55,above,font=\scriptsize] at (5.55,3.25) {汴河—都城水运轴};
  \draw[thick,dashed,color=timeline!85] (4.0,4.55) .. controls (5.55,4.5) and (6.25,3.95) .. (7.25,3.85);
  \node[text=timeline!85!black,above,font=\scriptsize] at (5.7,4.3) {陈桥至开封方向};

  \node[draw=timeline!80,fill=paper,rounded corners=2pt,align=center,inner sep=3pt] (chenqiao)
    at (3.0,2.55) {960年正月\\陈桥政权转换};
  \node[draw=dynasty!80,fill=paper,rounded corners=2pt,align=center,inner sep=3pt] (kaifeng)
    at (7.8,2.25) {承用后周中枢\\改国号宋};
  \draw[->,ultra thick,color=dynasty] (chenqiao.east) -- node[above,align=center,font=\scriptsize]
    {入据京师\\取得既有官僚与仓储} (kaifeng.west);

  \node[draw=source!70,fill=paper,rounded corners=2pt,align=center,inner sep=3pt] (supply)
    at (9.2,5.0) {政令、粮饷、兵员\\向东京汇集};
  \draw[->,thick,color=source] (supply.south west) -- (8.2,4.35);
  \node[text=source,align=center,font=\scriptsize] at (5.6,.85)
    {建国不只是一场仪式；都城、河运与禁军供给也须被接续};
\end{tikzpicture}

\smallskip
\small\textit{图\quad 陈桥—开封与北宋建国中枢示意：据官修纪传《宋史》卷1〈太祖本纪一〉、南宋编年汇集《续资治通鉴长编》建隆元年正月条、后世辑存的制度汇编《宋会要辑稿·礼》《宋会要辑稿·方域》，并参照开封北宋东京城遗址的城垣、街道与水系考古资料，说明政权转换与承用京师、仓储和水运节点的关系。黄袍仪式细节及都城空间承继层次均须谨慎。图中政权边界、水道和陈桥—开封路线均为约略示意。}
\end{center}


建国的日期和权力转换的基本过程，可由《宋史》太祖本纪与《续资治通鉴长编》建隆元年条相互勾连；建国礼仪则还须参看《宋会要辑稿·礼》。三种材料回答的不是同一个问题。《宋史》是元末官修史，给出后来王朝整理出的帝纪框架；《长编》由南宋李焘编年汇集宋廷旧材料，保存了更密的政事线索，却也带着后来的取舍；《宋会要辑稿》是自《永乐大典》辑出的制度条目，适合追问礼仪和机构，不能把它当作现场实录。后世最熟悉的“黄袍加身”图景，以及围绕陈桥的许多仪式性细节，正是在这些不同层次的叙述中逐渐定型。能够确定的是政权已经转换；不能据较晚的铺陈，补写帐中密语或人人共有的意图。

\begin{sourcenote}{陈桥能证实到哪里}
《宋史》《续资治通鉴长编》与《宋会要辑稿》可分别校读建国顺序、政事线索与礼仪语言；它们并非同一时代、同一体裁的现场记录。政权转换可以成立，军中密议和所有参与者的动机则不能由后出的铺陈补足。
\end{sourcenote}

宋初最急迫的功课，是使取得京城不再等于又一次禁军夺位。961年前后，太祖以任官、调动和优礼宿将等办法，削弱禁军统帅各自掌握的独立兵权。后世把它收束为一次饮宴中的“杯酒释兵权”，但《宋史》太祖本纪、石守信传与《长编》建隆二年条所能支持的，是一个持续的权力重组，而非一段可以逐字复原的对话。将领的个人声望被纳入任命与调动，军队的供养和指挥愈来愈回到东京；文官监督、轮调与财赋供给共同构成约束。它降低了五代式统兵者凭一支禁军改朝换代的风险，也使国家安全更依赖首都能够持续征粮、转运和支付。

\section{把割据地带接入中枢}

这种能力随后被用于统一，却不应把统一理解为一张地图在数年间自然合拢。963年宋军入荆南，结束高氏政权，取得长江中游的要点；《宋史》太祖本纪、《长编》乾德元年条和《宋会要辑稿·方域》的相关条目，让我们看到征服、改置州县和区域接收是相连的过程。荆南处在宋、后蜀和南唐之间，体量有限而位置关键；它的归入为沿江行动打开了较可靠的节点，并不意味着地方的旧有官员、运输和赋役立即换成同一种秩序。

964年冬伐后蜀，次年孟昶降宋，四川遂入宋境。成都平原的资源使西南具有吸引力，剑门道路和山地运输却提醒进军者：胜利后仍有粮道、驻军、降附者和旧官安置的问题。《宋史》的太祖本纪与后蜀世家提供了宋方的事件骨架，《长编》乾德二、三年条补足年序；军纪、赏赐的范围和地方社会的感受，则不能只由征服者的叙述断定。统一在这里同时是一项行政接收工程。后来蜀地的赋役与控制压力会以新的形式浮现，不能倒过来把965年的每一项安排都预言为必然失败。

岭南和江南的进入，进一步显示财政、交通与战争相互借力。971年宋军攻取广州，南汉刘鋹降服；宋朝由此接触岭南的港口、盐业和既有海上贸易网络，须在路州体系之外处理边地军政。974年渡江围金陵、975年李煜降宋，则把长江下游的攻城、漕运和江南财赋集中到同一场战争中。《宋史》南汉、南唐世家和太祖本纪同《长编》开宝四年、七年至八年条，支持这些节点的次序；战报里的兵数、伤亡和俘获数字却不能未经核对便化作确定的规模。金陵陷落后，宋廷获得江南的税粮、手工业和水网地区，但它面对的是接收而非空地：如何让旧政权辖下的官民、仓储与税源进入新朝的账目，才是统一得以延续的条件。

\phantomsection\label{fig:vol04:northern-song:early-state-capacity-transport}
\begin{center}
\begin{tikzpicture}[x=.88cm,y=.88cm,font=\small]
  \fill[paper] (0,0) rectangle (11.8,6.35);

  \fill[source!16] (.55,3.7) -- (3.55,3.5) -- (3.75,5.85) -- (.8,5.95) -- cycle;
  \node[align=center] at (2.12,4.8) {河东、河北\\边防与军需压力};
  \fill[timeline!21] (.7,.65) -- (4.0,.75) -- (3.7,2.9) -- (.5,3.0) -- cycle;
  \node[align=center] at (2.2,1.78) {四川\\山地道路、粮赋\\与接收后的治理};
  \fill[dynasty!19] (8.3,.65) -- (11.25,.8) -- (11.15,3.2) -- (8.1,2.95) -- cycle;
  \node[align=center] at (9.7,1.8) {江南\\水网、税粮、手工业\\与渡江战争};
  \fill[timeline!18] (8.1,3.65) -- (11.25,3.55) -- (11.3,5.9) -- (8.35,5.8) -- cycle;
  \node[align=center] at (9.85,4.85) {岭南\\港口、盐业与\\边地军政};

  \node[draw=dynasty!85,fill=dynasty!15,rounded corners=3pt,align=center,inner sep=5pt] (tokyo)
    at (5.85,3.15) {东京开封\\官署、仓储、禁军\\与文书调度};
  \node[draw=source!70,fill=paper,rounded corners=2pt,align=center,inner sep=3pt] (officials)
    at (5.85,5.18) {任命、轮调、文官监督\\与地方接收};
  \draw[<->,thick,color=source] (tokyo.north) -- node[right,font=\scriptsize]{调度、考核} (officials.south);

  \draw[->,ultra thick,color=source!85] (3.6,4.55) -- node[above,sloped,font=\scriptsize]{兵员、粮饷、信息} (5.15,3.65);
  \draw[->,ultra thick,color=timeline!90] (3.55,2.05) -- node[below,sloped,font=\scriptsize]{道路、赋役与接收} (5.12,2.75);
  \draw[->,ultra thick,color=dynasty!88] (8.05,1.95) -- node[below,sloped,font=\scriptsize]{水路、税粮与军需} (6.55,2.76);
  \draw[->,ultra thick,color=timeline!85] (8.05,4.65) -- node[above,sloped,font=\scriptsize]{港口、盐业与政务} (6.55,3.62);

  \node[text=source,align=center,font=\scriptsize] at (5.86,.75)
    {统一与守边依赖区域资源进入中枢；实际征收与运输并非单向无阻};
\end{tikzpicture}

\smallskip
\small\textit{图\quad 北宋初年区域资源、运输与国家能力关系示意：据官修纪传《宋史》卷1、卷2〈太祖本纪〉及后蜀、南汉、南唐世家，南宋编年汇集《续资治通鉴长编》乾德、开宝诸年条，后世辑存的制度汇编《宋会要辑稿·方域》，并参照开封北宋东京城遗址的水系考古资料，说明东京、边防、四川、江南和岭南之间的资源接收与调度关系。它不表示税额、运量、行政执行速度或某一地区必然服从中枢；图中区域边界、水道和资源、军需、信息路线均为约略示意。}
\end{center}


976年太祖去世，赵光义即位为太宗。即位日期可由《宋史》两朝本纪与《长编》开宝九年条确认；关于继承的“烛影斧声”却是后起传闻，现存材料并未使皇子与兄弟之间的安排完全透明。把这一不透明处改写成确定的宫廷阴谋，只会用后世兴趣覆盖当时的政治难题。太宗所继承的，是已扩大的财政与军队，也是尚未结束的河东割据和北方边患。

\section{河东之后，没有缓冲地带}

979年五月，宋灭北汉，河东归入宋朝，通常被视为五代十国割据大势的收束。《宋史》太宗本纪、北汉世家和《长编》太平兴国四年条足以确证灭国这一结果；北汉长期倚重辽的援助这一事实，也说明太原并非孤立城池。宋军取得河东后，宋与辽之间原先的缓冲政权随之消失。统一把军队和赋税推到了更远的北方，也把边界变成直接相接的战场与道路。

太宗旋即北攻幽州，在高梁河失利南撤。宋方的《宋史》《长编》和辽方的《辽史》景宗本纪都记下这场挫败，却在兵力、皇帝是否受伤、追击路线及功过归属上并不相同。因而高梁河不能只写成一方“失机”或另一方“天佑”的故事：宋军希望趁河东既下夺取燕云，辽军则能依托近畿、骑兵机动与既有防御应对。战败使速取燕云的设想落空，也迫使东京重新计算北部防线需要多少兵员、粮饷和指挥协调。

986年的雍熙北伐又以三路军北进，终至受挫撤军。三路并进的构想反映出宋廷仍想以协同动作牵制辽军；宋、辽两部史书与《长编》雍熙三年条同时提醒我们，诸路军的战果、军数和将领责任都经过战后叙事塑形。可以把握的后果是，北伐未能改变燕云的格局，守势边防和军费压力反而更深地进入中枢决策。此前限制将领独立兵权的制度，避免了五代式军头坐大，却也使远方战役更依赖来自东京的调度、文书和供给；这不是单一制度造成的胜败，却构成宋初军事协调所承受的约束。

997年赵恒继位为真宗，承接的正是这样的北边与财政。1004年冬，辽圣宗和萧太后南下，真宗至澶州前线，双方在军事对峙中转入议和。宋方的真宗本纪和《长编》景德元年冬条，与《辽史》圣宗本纪共同构成这一事件的基本证据；至于谁在战场上占有绝对主动、谁以何种言辞劝动皇帝，则多被后来的史家赋予戏剧性的形状。澶州的选择并不是抽象的怯战或勇战，而是两国都在机动、城防、补给和内部政治压力下衡量继续作战的成本。

1005年盟约落实，宋输送银绢，双方以兄弟之国交往。宋方记载中的“岁币”语言常被后世用来衡量屈辱或胜负；《宋会要辑稿·蕃夷》所存盟约条和《辽史》的并读，更适合把它放回两国关系的制度语境：银绢、使节、边市与防务被编入一套长期的停战安排，不能单凭一方的朝贡称谓推出单向统属。1008年的封禅和相关大礼，显示宋廷又尝试把边境安定转译为皇权与天命的语言；《宋史》真宗本纪、《长编》和《宋会要辑稿·礼》可见礼仪的举行与制度记录，却不能把祥瑞叙事当作事实的验证。

\phantomsection\label{fig:vol04:northern-song:song-liao-chanyuan-frontier}
\begin{center}
\begin{tikzpicture}[x=.88cm,y=.88cm,font=\small]
  \fill[paper] (0,0) rectangle (11.7,6.35);

  \fill[source!20] (.5,3.7) -- (11.2,3.45) -- (11.25,5.9) -- (.7,5.85) -- cycle;
  \node[align=center] at (9.5,4.95) {辽\\燕云、骑兵机动与北方近畿};
  \fill[dynasty!22] (.55,.55) -- (11.25,.7) -- (11.2,3.25) -- (.55,3.45) -- cycle;
  \node[align=center] at (2.0,1.72) {宋\\河北防务与东京中枢};

  \draw[very thick,color=timeline!85] (.7,3.52) .. controls (3.05,3.88) and (5.7,3.15) .. (8.0,3.55)
    .. controls (9.25,3.75) and (10.2,3.4) .. (11.15,3.55);
  \node[draw=timeline!80,fill=paper,rounded corners=2pt,align=center,inner sep=3pt] (line)
    at (5.95,3.6) {盟约维系的\\边境接触地带};

  \fill[timeline] (6.0,2.55) circle (.13);
  \node[below,align=center] at (6.0,2.45) {澶州\\1004年对峙};
  \fill[source] (3.4,4.6) circle (.13);
  \node[above,align=center] at (3.4,4.68) {燕云方向};
  \fill[dynasty] (9.55,1.65) circle (.13);
  \node[below,align=center] at (9.55,1.55) {东京\\开封};

  \draw[->,ultra thick,color=source!85] (3.65,4.45) .. controls (4.55,4.0) and (5.05,3.25) .. (5.75,2.68);
  \node[text=source,align=center,font=\scriptsize] at (4.2,3.4) {1004年辽军\\南下方向};
  \draw[->,ultra thick,color=dynasty] (9.3,1.72) .. controls (8.0,1.9) and (7.05,2.2) .. (6.25,2.52);
  \node[text=dynasty,align=center,font=\scriptsize] at (7.9,2.4) {宋廷赴澶州\\及河北动员};

  \draw[<->,thick,dashed,color=timeline!90] (5.45,3.06) -- (6.6,3.08);
  \node[text=timeline!90!black,above,align=center,font=\scriptsize] at (6.02,3.08)
    {使节、盟约与\\银绢输送关系};
  \node[draw=source!70,fill=paper,rounded corners=2pt,align=center,inner sep=3pt] at (9.3,4.8)
    {和平依赖\\防务、谈判、边市\\与地方管控并行};
  \node[text=source,align=center,font=\scriptsize] at (2.0,.85)
    {1005年后并非无边防的“固定国界”};
\end{tikzpicture}

\smallskip
\small\textit{图\quad 澶渊盟约前后的宋辽边境关系示意：据宋、辽官修纪传《宋史》卷5〈真宗本纪一〉与《辽史》〈圣宗本纪〉、南宋编年汇集《续资治通鉴长编》景德元年冬条，以及后世辑存的制度汇编《宋会要辑稿·蕃夷》辽国盟约条，呈现1004年澶州对峙、1005年盟约与防务、使节、银绢和边市之间的关联。双方战场主动性、银绢数额及各方称谓须回到文本语境辨析。图中政权边界、盟约接触地带与军队、使节和输送路线均为约略示意。}
\end{center}


盟约也没有把边界变成无人之地。1010年辽军再入宋境，宋廷加强河北边防，同时维持盟约渠道。《宋史》《辽史》与《长编》对破约和战果各有立场；能够看见的是，和平依赖的从来不只是纸上的誓词，还包括堡寨、侦察、地方管控、军队供给与使节往返。北宋最初数十年的国家能力，正在这种并行中成形：开封把资源向中枢汇集，统一战争把新的区域纳入账目与道路，北方压力又迫使朝廷不断校验它能调动多少人、粮和信息。它解决了五代以来最迫切的分裂与政变风险，却也把一套昂贵而持续的边防财政，留给了此后更长的宋辽关系。
